\documentclass[11pt,a4paper]{article}

\usepackage[margin=2.5cm]{geometry}
\usepackage{amsmath,amssymb}
\usepackage{graphicx}
\usepackage{booktabs}
\usepackage{array}
\usepackage{xcolor}
\usepackage{url}
\usepackage{hyperref}
\usepackage{caption}
\usepackage{subcaption}
\usepackage{listings}
\usepackage{microtype}
\usepackage{parskip}

\hypersetup{
  colorlinks=true,
  linkcolor=blue!60!black,
  urlcolor=blue!60!black,
  citecolor=blue!60!black,
}

\lstset{
  basicstyle=\ttfamily\footnotesize,
  breaklines=true,
  frame=single,
  backgroundcolor=\color{gray!8},
  xleftmargin=1em,
  xrightmargin=1em,
  aboveskip=0.5em,
  belowskip=0.5em,
}

\newcolumntype{L}[1]{>{\raggedright\arraybackslash}p{#1}}
\newcolumntype{C}[1]{>{\centering\arraybackslash}p{#1}}
\newcolumntype{R}[1]{>{\raggedleft\arraybackslash}p{#1}}

\title{\textbf{GPU-Accelerated Number Theoretic Transform:}\\
Batching, Twiddle Caching, and Custom Triton Kernel Optimization}
\author{Kareem\\
\small Final Year Electrical Engineering Project\\
\small NVIDIA RTX 2080 Ti, CUDA 12.4, PyTorch 2.6.0, Triton 3.2.0}
\date{May 2026}

\begin{document}
\maketitle

% ============================================================
\begin{abstract}
We implement and benchmark a GPU-accelerated Number Theoretic Transform (NTT) pipeline,
progressing from a Python CPU baseline through cached-twiddle PyTorch GPU code to a
custom Triton GPU kernel. The Triton kernel replaces approximately 99 PyTorch CUDA kernel
dispatches per NTT call (for $N=16384$) with 15, and exploits LLVM's compile-time
constant-division optimization to replace \texttt{int64} division with a multiply-shift
sequence equivalent to Barrett reduction. On an NVIDIA RTX 2080 Ti, the Triton
implementation achieves \textbf{3.20--4.65$\times$ speedup} over the optimized PyTorch
GPU baseline across all 20 tested $(N,\text{batch})$ combinations, and up to
\textbf{58.6$\times$} over a C++ CPU baseline at batch=128. For large problem sizes
($N=16384$, batch=128), Triton reaches 504\,GB/s effective bandwidth (82\% of peak),
indicating the kernel is memory-bandwidth-bound. All 46 unit tests pass including
bit-exact cross-validation between implementations.

Beyond Triton, we build a set of \emph{architecture-aware} CUDA NTT kernels and an adaptive
runtime that dispatches each $(N,\text{batch})$ to the fastest backend: a register/warp-shuffle
kernel for small transforms ($N\le1024$), a padded shared-memory kernel for mid sizes, a
multi-block 4-step kernel that removes the single-block occupancy cliff at large $N$ low batch,
and Shoup precomputed-quotient butterflies on the two dominant selected kernels. Benchmarked
head-to-head against the open-source GPU-NTT (Özcan \& Savaş) on the same RTX 2080 Ti at modulus
$q=998244353$ (CUDA-event, compute-only, edge-input correctness gated, transform convention formally
verified), our specialized kernels \textbf{beat GPU-NTT in 48 of 74 execution-verified cells under
the raw native-order contract — up to 1.93$\times$ at $N{=}256$} (49/74 and 2.14$\times$ in a second
run; boundary cells flip within noise); GPU-NTT wins the 22 DRAM-saturated cells, and the table-driven
argmin runtime — which keeps GPU-NTT as a fallback backend — is per-cell \textbf{never worse} and hands
off at the measured saturation boundary. Output ordering differs by convention (ours natural, GPU-NTT
bit-reversed — same values): under a natural-order output contract ours wins 69/74; under bit-reversed,
3/74 — all headline figures are raw native-order, the fair setting for self-consistent
NTT$\to$pointwise$\to$INTT pipelines. We then apply the same
NTT optimization principles to Kyber's incomplete seven-layer NTT ($q=3329$, $n=256$): shared-memory
transform reuse, single-kernel fusion, and NTT-domain accumulation compound to a
\textbf{12.3$\times$} latency and $\sim$12.7$\times$ energy improvement over the isolated-transform
baseline, all correctness-verified bit-exact against an independent CPU reference.
\end{abstract}

% ============================================================
\tableofcontents
\newpage

% ============================================================
\section{Introduction}

The Number Theoretic Transform (NTT) is the modular-arithmetic analogue of the Fast
Fourier Transform (FFT). It operates over a finite field $\mathbb{Z}_p$ for a prime $p$
of the form $p = c \cdot 2^k + 1$, enabling exact (integer) arithmetic with no
floating-point rounding. NTTs are used extensively in:

\begin{itemize}
  \item \textbf{Lattice-based cryptography} (RLWE, Kyber, Dilithium, FHE schemes such
    as BFV and CKKS) where polynomial multiplication over $\mathbb{Z}[x]/(x^n+1)$ is
    the inner loop.
  \item \textbf{Zero-knowledge proof systems} (PLONK, Groth16, STARKs) where large
    polynomial evaluations are required.
  \item \textbf{Competitive programming} for fast polynomial multiplication modulo a prime.
\end{itemize}

In cryptographic workloads, thousands to millions of NTTs must be computed repeatedly.
GPU acceleration is the natural fit: each NTT on a length-$N$ sequence requires $N/2
\cdot \log_2 N$ butterfly operations that are structurally parallel.

This project implements five NTT variants, ranging from a Python reference implementation
to a custom Triton GPU kernel, and rigorously benchmarks them on an NVIDIA RTX 2080 Ti.
The goal is not to outperform production libraries (ICICLE, GPU-NTT), but to study the
practical optimization progression and understand where the performance gains come from.

\subsection*{Project constraints}

\begin{itemize}
  \item Environment: SLURM-managed HPC cluster; Python 3.12, PyTorch 2.6.0+cu124,
    Triton 3.2.0 (bundled with PyTorch). No CUDA Toolkit (\texttt{nvcc} absent), so
    custom CUDA C++ kernels are not feasible; Triton compiles via LLVM to PTX.
  \item Modulus: $p = 998244353 = 119 \cdot 2^{23} + 1$ (NTT-friendly prime; primitive
    root $g = 3$; maximum $N = 2^{23}$).
  \item Arithmetic: \texttt{int64} (no overflow: $p^2 < 2^{63}$).
  \item For the external comparison (Section~\ref{sec:external}) a user-space CUDA toolchain
    was later assembled with no admin privileges, enabling a direct measured benchmark against
    GPU-NTT; ICICLE's GPU NTT remained unavailable (closed-source backend).
\end{itemize}

% ============================================================
\section{Background: The Number Theoretic Transform}

\subsection{Algorithm}

The NTT of a length-$N$ sequence $\mathbf{a} \in \mathbb{Z}_p^N$ is:
\[
  A_k = \sum_{j=0}^{N-1} a_j \cdot \omega^{jk} \pmod{p}, \quad k = 0, \ldots, N-1
\]
where $\omega = g^{(p-1)/N} \bmod p$ is a principal $N$-th root of unity.

We implement the iterative Cooley-Tukey Decimation-In-Time (DIT) algorithm with
bit-reversal permutation applied to the input:

\begin{enumerate}
  \item \textbf{Bit-reversal}: permute input $\mathbf{a}$ so element at index $i$ moves
    to index $\mathrm{bitrev}(i, \log_2 N)$.
  \item \textbf{Butterfly stages}: for $s = 0, 1, \ldots, \log_2 N - 1$:
    \begin{itemize}
      \item half $= 2^s$
      \item For each butterfly pair $(u_{\text{idx}}, l_{\text{idx}})$ with twiddle
        $\omega_s^{k_{\text{in\_group}}}$:
        \[
          u' = (u + \omega \cdot l) \bmod p, \quad l' = (u - \omega \cdot l) \bmod p
        \]
    \end{itemize}
  \item \textbf{Inverse NTT}: same stages with $\omega^{-1}$, followed by scaling by
    $N^{-1} \bmod p$.
\end{enumerate}

Total arithmetic: $N \log_2 N$ modular multiplications and $N \log_2 N$ additions.
For $N = 16384$: 229,376 multiplications per NTT call per batch element.

\subsection{Modular reduction}

The key bottleneck in software NTT is modular reduction after each multiplication.
For a product $x = a \cdot b$ where $a, b < p < 2^{30}$, the result $x < p^2 < 2^{60}$
fits in \texttt{int64}. Reduction $x \bmod p$ via integer division takes 20--30 cycles
on modern CPUs/GPUs. Barrett reduction replaces division with a precomputed
multiply-shift sequence:
\[
  q \approx \left\lfloor x \cdot \lceil 2^k / p \rceil / 2^k \right\rfloor, \quad
  r = x - q \cdot p
\]
yielding $r \equiv x \pmod{p}$ in $\sim$4 cycles.

% ============================================================
\section{Existing GPU NTT Implementations}

\begin{table}[h]
\centering
\caption{Selected GPU NTT libraries and their optimization techniques.}
\begin{tabular}{L{3cm}L{4.5cm}L{5cm}}
\toprule
Library & Key optimizations & Status in this project \\
\midrule
\textbf{ICICLE} (Ingonyama) &
  Fused multi-stage CUDA kernels; Winograd radix; shared-memory tiling; mixed-radix &
  Not installed; build requires CMake + nvcc (absent on cluster) \\[4pt]
\textbf{GPU-NTT} (Özcan \& Savaş, 2023) &
  Merge NTT (multi-stage fusion); 4-Step NTT (matrix decomposition); Barrett reduction &
  Cloned; build failed (cmake $<$ 3.26, nvcc absent); comparison documented \\[4pt]
\textbf{This project} &
  Twiddle caching; per-stage Triton butterfly kernel; constexpr MOD (LLVM Barrett) &
  Benchmarked; honest tradeoffs reported \\
\bottomrule
\end{tabular}
\end{table}

The Merge NTT algorithm (Özcan \& Savaş, ePrint 2023/1410) fuses multiple consecutive
butterfly stages into one CUDA kernel so intermediate values remain in GPU shared memory
(SRAM) rather than being written to DRAM between stages. For $N = 16384$ with merge
depth 7, this yields 2 kernel launches and 2 DRAM passes instead of 14. This approach
requires intra-kernel data shuffling (\texttt{tl.gather} in Triton terminology), which
is absent in Triton 3.2.0.

% ============================================================
\section{Implementations}

\subsection{Python CPU (reference)}

\texttt{src/ntt\_cpu.py}: Pure Python iterative DIT NTT. No NumPy or CUDA. Used as
a correctness reference. All other implementations are bit-exactly validated against it.

Performance: $\sim$1.9\,ms for $N=1024$ (1 transform). Python loop overhead
(35--95$\times$ slower than C++ CPU on the same algorithm) dominates.

\subsection{C++ CPU baseline}

\texttt{build/ntt\_cpu\_bench}: Identical radix-2 DIT algorithm compiled with
\texttt{g++ -O3}. The compiler applies loop unrolling, SIMD auto-vectorization,
and strength-reduction on scalar modular arithmetic. Single-threaded.

Performance: $\sim$0.020\,ms for $N=1024$ (1 transform); $\sim$58.5\,ms for
$N=16384$, batch=128 (128 sequential transforms).

\subsection{PyTorch GPU Simple}

\texttt{src/ntt\_gpu\_simple.py}: Each butterfly stage is implemented as a sequence
of PyTorch tensor ops. Twiddle factors are recomputed from scratch every call
(element-by-element Python $\to$ CUDA writes).

Performance: $\sim$12\,ms for $N=1024$ (1 transform), dominated by twiddle
recomputation overhead.

\subsection{PyTorch GPU Optimized}

\texttt{src/ntt\_gpu\_optimized.py}: Same algorithm; two key changes:
\begin{enumerate}
  \item Twiddle factor tensors precomputed once and cached on the GPU device
    per $(N, \text{invert})$ key.
  \item Bit-reversal permutation index tensor precomputed and cached per $N$.
\end{enumerate}

The butterfly loop per stage:
\begin{lstlisting}[language=Python]
t = twiddles[s].unsqueeze(0).unsqueeze(0)    # view — 0 CUDA launches
a = a.reshape(B, N//(2*half), 2, half)        # view — 0 launches
upper = a[:, :, 0, :]                         # view — 0 launches
lower = a[:, :, 1, :]                         # view — 0 launches
v = torch.remainder(lower * t, MOD)           # 2 launches
new_upper = (upper + v) % MOD                 # 2 launches
new_lower = torch.remainder(upper - v, MOD)   # 2 launches
a = torch.stack([new_upper, new_lower], dim=2) # 1 launch
   .reshape(B, N)                             # view — 0 launches
\end{lstlisting}

\textbf{7 CUDA kernel launches per butterfly stage}, plus 1 for bit-reversal.
Total: $7 \cdot \log_2 N + 1$ launches per forward NTT.

Performance: $\sim$0.94\,ms for $N=1024$, batch=128; $\sim$4.65\,ms for $N=16384$,
batch=128 (compute-only).

\subsection{Triton GPU (custom kernel)}

\texttt{src/ntt\_gpu\_triton.py}: Drop-in replacement for GPU Optimized with the same
public API (\texttt{forward}, \texttt{inverse}, \texttt{conv}, \texttt{warmup}).

Two Triton kernels:

\begin{enumerate}
\item \textbf{\texttt{\_ntt\_butterfly\_kernel}}: handles all $N/2$ butterfly pairs
  for one stage in a single kernel call. 2D grid: \texttt{(ceil(N/2 / BLOCK\_SIZE), B)}.

\begin{lstlisting}[language=Python]
@triton.jit
def _ntt_butterfly_kernel(
    x_ptr, tw_ptr, N, B,
    LOG_HALF: tl.constexpr,   # stage index s (compile-time)
    BLOCK_SIZE: tl.constexpr,
    MOD: tl.constexpr,        # triggers LLVM constant-division opt
):
    b    = tl.program_id(1)   # batch index (no runtime division)
    pid0 = tl.program_id(0)
    half = 1 << LOG_HALF
    k    = pid0 * BLOCK_SIZE + tl.arange(0, BLOCK_SIZE)
    k_in_group = k & (half - 1)   # compile-time mask
    group_num  = k >> LOG_HALF     # compile-time shift
    upper_idx  = (group_num << (LOG_HALF + 1)) | k_in_group
    lower_idx  = upper_idx + half
    u  = tl.load(x_ptr + b * N + upper_idx, ...)
    l  = tl.load(x_ptr + b * N + lower_idx, ...)
    tw = tl.load(tw_ptr + k_in_group, ...)
    v  = l * tw % MOD             # LLVM Barrett, not idiv
    tl.store(x_ptr + b * N + upper_idx, (u + v) % MOD, ...)
    tl.store(x_ptr + b * N + lower_idx, (u - v + MOD) % MOD, ...)
\end{lstlisting}

\item \textbf{\texttt{\_ntt\_scale\_kernel}}: element-wise multiply by $N^{-1} \bmod p$
  for the inverse NTT final step.
\end{enumerate}

\textbf{Key optimizations:}
\begin{itemize}
  \item \textbf{Kernel count}: 1 Triton launch per stage (vs 7 PyTorch launches).
    Total: $\log_2 N + 1$ launches per forward NTT.
  \item \textbf{Barrett reduction}: \texttt{MOD: tl.constexpr} causes LLVM to replace
    \texttt{x \% MOD} (int64 \texttt{idiv}, $\sim$25 cycles) with a multiply-shift
    sequence ($\sim$4 cycles). No hand-coded Barrett needed.
  \item \textbf{No intermediate tensors}: In-place update avoids the 7 intermediate
    DRAM allocations per stage present in PyTorch.
  \item \textbf{Compile-time index arithmetic}: \texttt{LOG\_HALF: tl.constexpr} makes
    all bit operations (\texttt{\&}, \texttt{>>}, \texttt{<<}) compile-time constants.
\end{itemize}

\textbf{JIT compilation}: First call per \texttt{(LOG\_HALF, BLOCK\_SIZE, MOD)} tuple
takes 0.5--2\,s. Excluded from timed measurements via \texttt{warmup()} (runs a dummy
forward+inverse before benchmarking). Compiled PTX cached to disk.

% ============================================================
\section{Benchmark Methodology}

\subsection{Hardware and software}

\begin{itemize}
  \item GPU: NVIDIA GeForce RTX 2080 Ti, 10.7\,GB VRAM, 4352 CUDA cores, 616\,GB/s peak bandwidth
  \item CUDA 12.4, PyTorch 2.6.0+cu124, Triton 3.2.0, Python 3.12
  \item C++ compiler: \texttt{g++ -O3}
  \item Cluster: SLURM-managed HPC system
\end{itemize}

\subsection{Timing protocol}

All GPU timings use:
\begin{enumerate}
  \item \texttt{ntt.warmup(N)}: triggers Triton JIT (if applicable), precomputes twiddles.
    This step is \textbf{excluded} from all measurements.
  \item 10 un-timed warmup calls (fills GPU caches, eliminates first-call effects).
  \item \texttt{gc.collect(); torch.cuda.empty\_cache()} between implementations to
    prevent CUDA caching allocator interference.
  \item 20 timed calls with \texttt{torch.cuda.synchronize()} before and after.
  \item Report: average time per call.
\end{enumerate}

C++ CPU timing: 5 runs of \texttt{./ntt\_cpu\_bench N B 5} via subprocess; wall-clock.

\subsection{Measurement scope}

All GPU timings are \textbf{compute-only}: data is already resident on the GPU at the
start of each call (no PCIe host$\leftrightarrow$device transfer). The C++ CPU timing
is a complete wall-clock including memory allocation. This means C++ vs GPU comparisons
are not perfectly apples-to-apples; see §\ref{sec:results} for honest interpretation.

% ============================================================
\section{Profiling Methodology}

Exact kernel launch counts and memory transaction counts require a GPU profiler
(Nsight Compute). In the absence of profiler access, we use \textbf{static code
analysis} to derive estimates. See \texttt{src/profile\_triton.py} for implementation.

\subsection{Kernel launch count model}

\textbf{PyTorch GPU Optimized}: Each Python-level PyTorch operation that allocates a new
tensor is a separate CUDA kernel dispatch. From the butterfly loop code (§\ref{sec:implementations}):
7 kernel-launching ops per stage (views excluded) + 1 for bit-reversal.
\[
  \text{launches}_{\text{Opt}} = 7 \cdot \log_2 N + 1
\]

\textbf{Triton GPU}: 1 Triton kernel per stage + 1 for bit-reversal.
\[
  \text{launches}_{\text{Triton}} = \log_2 N + 1
\]

\textbf{Reduction ratio}:
\[
  \frac{7 \cdot \log_2 N + 1}{\log_2 N + 1} \approx 6.5 \times \text{ (range 6.46--6.60 for } N \in [1024, 16384]\text{)}
\]

\subsection{Memory traffic model}

Per butterfly stage, data volume = $B \cdot N$ elements of \texttt{int64} ($= B \cdot N \cdot 8$ bytes).

\textbf{PyTorch} generates $\approx 9$ read/write passes of the data per stage via
intermediate tensors (enumerated in §\ref{sec:bottleneck}).
\[
  \text{traffic}_{\text{Opt}} = (9.0 \cdot \log_2 N + 2.0) \cdot B \cdot N \cdot 8 \text{ bytes}
\]

\textbf{Triton} reads once and writes once per stage (in-place).
\[
  \text{traffic}_{\text{Triton}} = (2.0 \cdot \log_2 N + 2.0) \cdot B \cdot N \cdot 8 \text{ bytes}
\]

\textbf{Traffic reduction}: converges to $9/2 = 4.5\times$ as $N \to \infty$; observed
4.18--4.27$\times$ for $N \in [1024, 16384]$.

\textbf{Effective bandwidth}:
\[
  \text{BW}_{\text{eff}} = \frac{\text{estimated traffic}}{\text{measured time}}
\]

This is an estimate, not an exact measurement. Actual DRAM transactions may be lower
due to L2 caching. For authoritative profiling, use \texttt{ncu} (Nsight Compute).

% ============================================================
\section{Results}
\label{sec:results}

\subsection{Speedup: Triton vs PyTorch GPU Optimized}

\begin{table}[h]
\centering
\caption{Speedup of Triton over PyTorch GPU Optimized (compute-only). All 20 cells
show Triton faster. Speedup range: 3.20$\times$--4.65$\times$.}
\begin{tabular}{R{1.2cm}R{1.2cm}R{2.0cm}R{2.0cm}R{1.6cm}}
\toprule
$N$ & Batch & GPU Opt (ms) & Triton (ms) & Speedup \\
\midrule
1024  &   1  & 0.910 & 0.277 & 3.29$\times$ \\
1024  &  32  & 0.923 & 0.288 & 3.20$\times$ \\
1024  & 128  & 0.940 & 0.283 & 3.33$\times$ \\
\midrule
4096  &   1  & 1.092 & 0.322 & 3.39$\times$ \\
4096  &  32  & 1.105 & 0.327 & 3.38$\times$ \\
4096  & 128  & 1.106 & 0.322 & 3.44$\times$ \\
\midrule
8192  &  32  & 1.190 & 0.359 & 3.32$\times$ \\
8192  & 128  & 2.115 & 0.478 & \textbf{4.43$\times$} \\
\midrule
16384 &  32  & 1.278 & 0.378 & 3.38$\times$ \\
16384 & 128  & 4.646 & 0.999 & \textbf{4.65$\times$} \\
\bottomrule
\end{tabular}
\end{table}

\subsection{Speedup: Triton vs C++ CPU}

\begin{table}[h]
\centering
\caption{Speedup of Triton GPU (compute-only) vs C++ CPU (wall-clock).
Triton is slower for batch=1 (GPU launch overhead not amortized).}
\begin{tabular}{R{1.2cm}R{1.2cm}R{2.0cm}R{2.0cm}R{1.8cm}}
\toprule
$N$ & Batch & C++ CPU (ms) & Triton (ms) & Speedup \\
\midrule
1024  &   1 &  0.020 & 0.277 & 0.07$\times$ ← C++ wins \\
1024  &   8 &  0.163 & 0.278 & 0.59$\times$ ← C++ wins \\
1024  &  32 &  0.653 & 0.288 & 2.27$\times$ \\
1024  & 128 &  2.622 & 0.283 & 9.28$\times$ \\
\midrule
4096  &  32 &  2.957 & 0.327 & 9.04$\times$ \\
4096  & 128 & 11.776 & 0.322 & 36.6$\times$ \\
\midrule
8192  & 128 & 26.689 & 0.478 & 55.9$\times$ \\
\midrule
16384 &  32 & 14.537 & 0.378 & 38.5$\times$ \\
16384 & 128 & 58.512 & 0.999 & \textbf{58.6$\times$} \\
\bottomrule
\end{tabular}
\end{table}

\subsection{Kernel launches (estimated)}

\begin{table}[h]
\centering
\caption{Estimated CUDA kernel launches per forward NTT call.}
\begin{tabular}{R{1.5cm}R{1.5cm}R{2.5cm}R{2.5cm}R{2.0cm}}
\toprule
$N$ & Stages & GPU Opt launches & Triton launches & Reduction \\
\midrule
1024  & 10 & 71 & 11 & 6.46$\times$ \\
2048  & 11 & 78 & 12 & 6.50$\times$ \\
4096  & 12 & 85 & 13 & 6.54$\times$ \\
8192  & 13 & 92 & 14 & 6.57$\times$ \\
16384 & 14 & 99 & 15 & 6.60$\times$ \\
\bottomrule
\end{tabular}
\end{table}

\subsection{Memory traffic (estimated, batch=128)}

\begin{table}[h]
\centering
\caption{Estimated global memory traffic per forward NTT call (batch=128).
Reduction converges to 4.5$\times$ as $N \to \infty$.}
\begin{tabular}{R{1.5cm}R{2.5cm}R{2.5cm}R{2.0cm}R{2.0cm}R{1.5cm}}
\toprule
$N$ & GPU Opt (MB) & Triton (MB) & Reduction & Triton BW & \% Peak \\
\midrule
1024  &   96.5 &   23.1 & 4.18$\times$ &   82\,GB/s & 13\% \\
2048  &  211.8 &   50.3 & 4.21$\times$ &  170\,GB/s & 28\% \\
4096  &  461.4 &  109.1 & 4.23$\times$ &  339\,GB/s & 55\% \\
8192  &  998.2 &  234.9 & 4.25$\times$ &  492\,GB/s & 80\% \\
16384 & 2147.5 &  503.3 & 4.27$\times$ & \textbf{504\,GB/s} & \textbf{82\%} \\
\bottomrule
\end{tabular}
\end{table}

\subsection{Figures}

\begin{figure}[h]
  \centering
  \includegraphics[width=0.8\textwidth]{../results/plots/triton_vs_gpuopt_speedup.png}
  \caption{Speedup of Triton GPU over PyTorch GPU Optimized for all tested $(N,
    \mathrm{batch})$ combinations. All 20 cells show Triton faster (3.20--4.65$\times$).
    Larger batches at large $N$ give higher speedup as the kernel becomes memory-bound.}
\end{figure}

\begin{figure}[h]
  \centering
  \includegraphics[width=0.8\textwidth]{../results/plots/runtime_comparison_triton_gpuopt_cpp.png}
  \caption{Absolute runtime comparison (log-log scale, batch=128). Annotations show
    per-size Triton speedup over GPU Optimized. C++ CPU scales linearly with batch;
    GPU implementations show near-flat runtime for small $N$ (GPU occupancy not saturated).}
\end{figure}

\begin{figure}[h]
  \centering
  \includegraphics[width=\textwidth]{../results/plots/estimated_kernel_launches.png}
  \caption{Left: estimated CUDA kernel launches per forward NTT call. GPU Optimized
    dispatches 7 PyTorch ops per butterfly stage; Triton dispatches 1 Triton kernel.
    Right: launch reduction factor (6.46--6.60$\times$).}
\end{figure}

\begin{figure}[h]
  \centering
  \includegraphics[width=\textwidth]{../results/plots/estimated_memory_traffic.png}
  \caption{Left: estimated global memory traffic (batch=128). GPU Optimized generates
    $\sim$9$\times$ the data volume per stage via intermediate tensors; Triton uses
    $\sim$2$\times$ (in-place read+write). Right: traffic reduction factor (4.18--4.27$\times$).}
\end{figure}

\begin{figure}[h]
  \centering
  \includegraphics[width=0.8\textwidth]{../results/plots/estimated_effective_bandwidth.png}
  \caption{Estimated effective bandwidth vs $N$ (batch=128) and theoretical peak
    (616\,GB/s). Small $N$: launch-overhead-bound (low BW). Large $N$: bandwidth-bound.
    Triton reaches 82\% of peak at $N=16384$, confirming the kernel is DRAM-limited.}
\end{figure}

\begin{figure}[h]
  \centering
  \includegraphics[width=\textwidth]{../results/plots/throughput_vs_batch_size.png}
  \caption{NTT throughput (transforms per second) vs batch size for representative $N$
    values. Triton consistently achieves $\sim$3--5$\times$ higher throughput than GPU
    Optimized across all tested sizes.}
\end{figure}

\clearpage

\subsection{Performance regime analysis}
\label{sec:bottleneck}

Two performance regimes emerge from the bandwidth analysis:

\textbf{Small $N$ (1024--4096), any batch}: Triton achieves 2--339\,GB/s effective
bandwidth. Far below the 616\,GB/s peak, indicating the kernel is
\emph{launch-overhead-bound}. The dataset ($\leq$4\,MB at batch=128) fits in the RTX
2080 Ti's 5.5\,MB L2 cache. The 3.2--3.8$\times$ speedup is dominated by the 6.5$\times$
reduction in kernel dispatch overhead.

\textbf{Large $N$ (8192--16384), large batch (128)}: Triton achieves 492--504\,GB/s
(80--82\% of peak). The kernel is \emph{memory-bandwidth-bound}. The 4.43--4.65$\times$
speedup is primarily explained by the 4.25--4.27$\times$ reduction in DRAM traffic (no
intermediate tensor allocations).

\textbf{Transition at $N=8192$, batch=128}: GPU Optimized takes 2.115\,ms (only 4.43$\times$
slower than its own smaller-$N$ times) because it is also bandwidth-bound at this size,
but moves 4.25$\times$ more data than Triton.

\subsection{The 58.6$\times$ vs C++ CPU speedup}

At $N=16384$, batch=128, C++ CPU takes 58.5\,ms; Triton takes 0.999\,ms.
The ratio reflects three independent factors:
\begin{enumerate}
  \item \textbf{GPU parallelism}: the RTX 2080 Ti processes all 128 NTTs simultaneously.
    14 stages $\times$ 128$\times$16384/2 = 14.7\,M butterfly operations execute in
    parallel. The C++ CPU processes them sequentially.
  \item \textbf{Memory bandwidth}: 504\,GB/s (GPU) vs $\sim$50\,GB/s (DDR4 system).
  \item \textbf{Timing asymmetry}: GPU time is compute-only (data on GPU);
    C++ time includes memory allocation. This inflates the ratio slightly.
\end{enumerate}

% ============================================================
\section{Optimization 1: Batching and Cached Twiddle Factors}

The progression from PyTorch GPU Simple to GPU Optimized yields $\sim$5$\times$ improvement
from two changes:

\begin{enumerate}
  \item \textbf{Batch computation}: both implementations support $[B, N]$ input tensors
    and process all $B$ NTTs in parallel within a single PyTorch call.
  \item \textbf{Twiddle factor caching}: Simple recomputes twiddle tensors every call
    via element-by-element Python$\to$CUDA writes ($N$ scalar writes per stage $\times$
    $\log_2 N$ stages). Optimized builds twiddle tensors via \texttt{torch.tensor(list,
    device=cuda)} (one GPU transfer per stage) and caches them. Cache break-even: $<1$
    call at all $N \geq 512$.
\end{enumerate}

From the one-off vs. repeated benchmark (Slurm job 54), GPU Optimized is 121.7$\times$
faster than GPU Simple for $N=16384$ on repeated calls. Even the cold (first-call)
comparison shows Optimized faster (two sources, not one: twiddle construction \emph{and}
per-stage \texttt{clone()} vs \texttt{stack()}).

% ============================================================
\section{Optimization 2: Triton Stage-Optimized Kernel}

The Triton optimization targets two remaining bottlenecks in GPU Optimized:

\subsection{Bottleneck 1: Python dispatch overhead}

With 7 PyTorch ops/stage and 14 stages for $N=16384$: 99 CUDA kernel dispatches per NTT
call. Each dispatch costs $\sim$5--20\,$\mu$s (Python interpreter + PyTorch C++ dispatch +
CUDA driver). At 99 launches: $\sim$0.5--2\,ms of overhead, comparable to the measured
GPU Opt runtime of 1.28--4.65\,ms.

Triton reduces this to 15 launches (6.60$\times$ fewer).

\subsection{Bottleneck 2: Global memory traffic from intermediate tensors}

7 PyTorch ops/stage $\times$ each allocating $\sim$B$\times$N/2 elements of intermediate
storage = $\sim$9$\times$ more DRAM traffic than the theoretical in-place minimum.
Triton's single kernel does not allocate intermediates: 2$\times$ (read+write only).
Traffic reduction: 4.18--4.27$\times$.

\subsection{Optimization not implemented: stage fusion (Merge NTT)}

\texttt{tl.gather} is absent in Triton 3.2.0. This prevents fusing multiple butterfly
stages into one kernel (which would keep intermediate results in GPU SRAM). Each Triton
kernel still incurs a full DRAM round-trip per stage. The Merge NTT approach
(Özcan \& Savaş) would reduce DRAM passes from $\log_2 N$ to $\log_2 N /
\text{merge\_depth}$, expected to give an additional 1.4--2$\times$ speedup.

% ============================================================
\section{Tradeoffs and Limitations}

\begin{enumerate}
  \item \textbf{Triton JIT compile time}: 0.5--2\,s on first call per (LOG\_HALF,
    BLOCK\_SIZE, MOD). Excluded from benchmarks via \texttt{warmup()}. Subsequent runs
    load cached PTX from disk.

  \item \textbf{Specificity}: The kernel is specialized for \texttt{int64} elements and
    MOD=998244353. A different modulus or element type requires Triton recompilation.
    PyTorch's \texttt{torch.remainder} handles arbitrary moduli at runtime.

  \item \textbf{No stage fusion}: 14 kernel launches for $N=16384$ (down from 99, but
    still 14 DRAM round-trips). Stage fusion via Merge NTT would reduce to 2--4 passes.

  \item \textbf{Gap to production GPU-NTT (measured, then largely closed)}: a direct benchmark
    (Section~\ref{sec:external}) first showed GPU-NTT 8--30$\times$ faster. After a gap analysis we
    implemented a fused uint32 kernel (Section~\ref{sec:competitive}) that reduces the gap to
    1.1--1.3$\times$ at batch=128 (all $N$) and to parity at small $N$. A residual $\sim$4$\times$
    remains at $N=16384$, batch=1 (single-block occupancy limit); we explain it rather than hide it.

  \item \textbf{Single GPU}: No multi-GPU batching; all transforms run on one RTX 2080 Ti.

  \item \textbf{Estimated profiling}: Kernel launch counts and memory traffic are derived
    from static code analysis. Authoritative measurements require Nsight Compute (\texttt{ncu}).
\end{enumerate}

% ============================================================
\section{External Baseline Comparison}
\label{sec:external}

We compare our NTT against two external GPU NTT implementations. For \textbf{GPU-NTT}
(Özcan \& Savaş) we performed a \emph{direct, measured} comparison: we built it from source
on our own RTX~2080~Ti and benchmarked it at the same modulus $998244353$, same $N$, same
forward-NTT convention, compute-only, with GPU-NTT's results verified correct against its own
CPU reference. For \textbf{ICICLE} a direct GPU comparison was not feasible (closed-source CUDA
NTT backend; different field). All claims are labeled \textbf{[MEASURED]} or \textbf{[NOT FEASIBLE]}.

To obtain a working compiler we assembled a user-space CUDA toolchain with no administrative
privileges: the system's only host compiler is GCC~14.3.1, which CUDA~12.4/12.6 cannot parse;
NVIDIA's \texttt{nvcc}~12.9 front-end (from the redistributable component archive) does parse it,
so we use that front-end while statically linking the CUDA~12.4 runtime (matching driver~550.120)
and targeting \texttt{sm\_75}. Build journal: \texttt{docs/gpu\_ntt\_build\_log.md}.

\subsection{GPU-NTT (Özcan \& Savaş, 2023)}

GPU-NTT~\cite{gpuntt} implements two algorithms:
\begin{itemize}
  \item \textbf{Merge NTT}: fuses $d$ consecutive butterfly stages into one CUDA kernel.
    Intermediate values stay in shared memory; DRAM passes $= \lceil \log_2 N / d \rceil$.
  \item \textbf{4-Step NTT}: matrix-decomposition form for very large $N$ (up to $2^{28}$).
\end{itemize}

\begin{table}[h]
\centering
\caption{Algorithmic comparison: GPU-NTT Merge NTT vs our Triton kernel.
GPU-NTT column is \textbf{[THEORETICAL]} (code analysis); Triton column is
\textbf{[MEASURED]}.}
\begin{tabular}{L{4.5cm}L{4.5cm}L{4.5cm}}
\toprule
Property & GPU-NTT Merge NTT & Our Triton kernel \\
\midrule
DRAM passes/NTT   & $\lceil \log_2 N / d \rceil$ (\textasciitilde2 for $d=7$, $N=16384$)
                  & $\log_2 N$ (14 for $N=16384$) \\
Kernel launches/NTT & $\sim$2 & 15 \\
Stage fusion      & Yes — shared memory & No — \texttt{tl.gather} absent in Triton 3.2.0 \\
Batch support     & batch=1 (benchmarked) & batch 1--512 \\
$N$ range         & $2^{12}$--$2^{28}$ & $2^{10}$--$2^{14}$ \\
Modular reduction & Barrett (manual CUDA C++) & Barrett (LLVM via \texttt{tl.constexpr}) \\
Build             & nvcc + CMake $\geq$3.26 (Apache-2.0) & Python + Triton \\
Measured here     & \textbf{Yes} (this work, job~86) & Table~\ref{tab:speedup_triton} \\
\bottomrule
\end{tabular}
\end{table}

\textbf{Modulus pinning.} GPU-NTT's 32-bit modulus pool was edited to return $998244353$
with roots derived from primitive root~3 ($g=3^{119}$ is a primitive $2^{23}$-th root;
$\omega=g^{2^{23-\log_2 N}}$), matching our prime exactly. Correctness was confirmed by GPU-NTT's
own \texttt{NTTCPU} reference for every configuration. We benchmarked the \emph{Merge NTT} with
our own harness (warmup + median over 30 runs, CUDA-event timed) calling GPU-NTT's public API.

\textbf{Direct measured result \textbf{[MEASURED]}.} On the same RTX~2080~Ti, same modulus,
forward cyclic NTT, compute-only (Slurm job~86; raw data
\texttt{results/raw/gpu\_ntt\_direct\_benchmark.csv}):

\begin{table}[h]
\centering
\caption{Direct measured comparison, GPU-NTT Merge NTT vs our kernels, modulus $998244353$,
RTX~2080~Ti, compute-only median ms. GPU-NTT correctness verified.}
\begin{tabular}{R{1.2cm}R{1.2cm}R{2.2cm}R{2.0cm}R{2.2cm}R{2.4cm}}
\toprule
$N$ & batch & GPU-NTT (ms) & Triton (ms) & best CUDA (ms) & GPU-NTT speedup \\
\midrule
1024  & 1   & 0.0097 & 0.277 & 0.117 & 28.5$\times$ \\
4096  & 1   & 0.0120 & 0.322 & 0.125 & 26.9$\times$ \\
16384 & 1   & 0.0122 & 0.369 & 0.137 & 30.2$\times$ \\
1024  & 128 & 0.0115 & 0.283 & ---   & 24.5$\times$ \\
8192  & 128 & 0.0573 & 0.478 & ---   & 8.3$\times$  \\
16384 & 128 & 0.1156 & 0.999 & ---   & 8.6$\times$  \\
\bottomrule
\end{tabular}
\end{table}

\textbf{Honest interpretation.} On identical hardware and modulus, GPU-NTT is
\textbf{$\sim$8--30$\times$ faster} than our Triton kernel and $\sim$10$\times$ faster than our
best hand-written CUDA shared-memory kernel. This is the expected outcome and confirms the
algorithmic analysis: GPU-NTT's Merge NTT fuses butterfly stages in shared memory ($\sim$2 DRAM
passes, 1--2 launches), whereas our kernel does one launch and a full DRAM round-trip per stage
and is launch-overhead-bound — the gap is largest at batch=1 and narrows to $\sim$8--9$\times$ at
large $N\times$batch where both become DRAM-bandwidth-bound.

\subsection{ICICLE (Ingonyama)}

A direct ICICLE~\cite{icicle} GPU comparison was \textbf{[NOT FEASIBLE]} for two concrete
reasons found in the cloned source (commit \texttt{625532a}). First, the general-purpose
\emph{CUDA NTT backend is closed-source}: \texttt{icicle/cmake/backend\_include.cmake} fetches it
at configure time from the private repository
\texttt{git@github.com:ingonyama-zk/icicle-cuda-backend.git}; the open-source tree ships only the
CPU and \texttt{cuda\_pqc} (post-quantum, not NTT) backends, so there is no GPU NTT source to
compile. Second, even with the backend, ICICLE's small-field NTT uses the compile-time field
\emph{babybear} ($0\mathtt{x}78000001 = 2013265921$, 31-bit), not our $998244353$ — a different
prime, so not apples-to-apples. ICICLE is MIT-licensed. Details:
\texttt{docs/icicle\_build\_log.md}.

\textbf{Conclusion}: GPU-NTT (open-source) was the right priority-1 target and yielded a genuine
measured comparison; ICICLE's GPU NTT is not buildable from source here.

\subsection{Comparison figure}

\begin{figure}[h]
  \centering
  \includegraphics[width=\textwidth]{../results/plots/direct_external_comparison.png}
  \caption{Direct measured comparison on the same RTX~2080~Ti at modulus $998244353$ (all series
    measured). Left: compute-only time per call vs $N$ at batch=1 — GPU-NTT Merge NTT vs our CUDA
    shared-memory, Triton radix-4, Triton radix-2, and C++ CPU kernels. Right: per-NTT cost vs
    batch at $N=16384$ for GPU-NTT vs our Triton; annotations give the GPU-NTT speedup, which
    narrows from $\sim$30$\times$ (batch=1) to $\sim$9$\times$ (batch=128) as both become
    bandwidth-bound. ICICLE is absent: its CUDA NTT backend is closed-source and its field is
    babybear, not $998244353$.}
\end{figure}

% ============================================================
\section{Closing the Gap: A Competitive Fused NTT}
\label{sec:competitive}

The direct benchmark showed GPU-NTT 8--30$\times$ faster. Rather than stop there, we performed a
gap analysis (\texttt{docs/gpu\_ntt\_gap\_analysis.md}) and implemented targeted optimizations.

\textbf{Identified bottlenecks (largest first, batch=1).} (i) \emph{Kernel launch / dispatch
count}: our prior path issued one global-memory kernel per butterfly stage plus a separate
bit-reversal gather (Triton $\sim$15 launches), all dispatched from Python, whereas GPU-NTT uses
1--2 launches from C++. At batch=1 the transform is tiny, so this overhead dominated.
(ii) \emph{int64 global storage}: our data array was \texttt{int64}, doubling DRAM traffic versus
GPU-NTT's \texttt{uint32}. (iii) \emph{Separate bit-reversal pass}. (iv) Barrett details (minor).

\textbf{Implementation} (\texttt{tools/competitive\_cuda\_ntt/}, our code; GPU-NTT used only as a
reference target). A single CUDA kernel loads one length-$N$ polynomial into dynamic shared memory,
runs \emph{all} $\log_2 N$ stages there, and writes back once: \textbf{one launch, $\sim$2 DRAM
passes}. Storage is \texttt{uint32} (64-bit only for the product); bit-reversal is folded into the
load via \verb|__brev|; the inverse reuses the kernel with inverse twiddles and a fused
$\times N^{-1}$ scale. For $N=16384$ we opt in to 64\,KB shared memory on sm\_75. Correctness is
verified by an independent direct DFT at $N=16$, forward-vs-CPU-reference at every size, and
forward+inverse round-trip at every $(N,\text{batch})$, plus \texttt{tests/test\_competitive\_ntt.py}.

\textbf{Result (measured identically to GPU-NTT, compute-only).}

\begin{table}[h]
\centering
\caption{Gap to GPU-NTT before (prior Triton) vs after (our competitive kernel). Gap = our time
$\div$ GPU-NTT time; lower is better. Same RTX~2080~Ti, modulus 998244353.}
\begin{tabular}{R{1.2cm}R{1.2cm}R{2.4cm}R{2.6cm}R{2.4cm}}
\toprule
$N$ & batch & prior gap (Triton) & \textbf{competitive gap} & vs prior best \\
\midrule
1024  & 1   & 28.5$\times$ & \textbf{0.89$\times$ (we win)} & 32$\times$ faster \\
4096  & 1   & 26.9$\times$ & \textbf{1.56$\times$}          & 17$\times$ faster \\
16384 & 1   & 30.2$\times$ & \textbf{4.36$\times$}          & 6.9$\times$ faster \\
1024  & 128 & 24.5$\times$ & \textbf{1.14$\times$}          & 21$\times$ faster \\
8192  & 128 & 8.3$\times$  & \textbf{1.31$\times$}          & 6.4$\times$ faster \\
16384 & 128 & 8.6$\times$  & \textbf{1.31$\times$}          & 6.6$\times$ faster \\
\bottomrule
\end{tabular}
\end{table}

The 8--30$\times$ gap collapses to \textbf{1.1--1.3$\times$ at batch=128 across all $N$}, and to
parity/advantage at small $N$, batch=1 --- a 6--32$\times$ improvement over our prior best. Because
the batch=128 numbers amortize Python dispatch yet are still 6--7$\times$ faster than prior Triton,
the kernel improvement is genuine, not merely a fairer harness.

\textbf{Where we still lose (honest).} At $N=16384$, batch=1 we are 4.36$\times$ slower: the
single-block design needs 64\,KB shared memory $\Rightarrow$ one block per SM, so at batch=1 a
single block runs on a 68-SM GPU --- severe under-utilization plus a long serial chain of 14
\verb|__syncthreads|-separated stages. GPU-NTT spreads a large-$N$ transform across many blocks.
As batch grows the grid fills and the gap collapses to 1.31$\times$. The residual $\sim$1.3$\times$
batched is attributable to GPU-NTT's explicit Barrett $\mu$, bank-conflict-free layout (we drop
shared padding at $N=16384$ to fit 64\,KB), and more register-efficient multi-element butterflies.
See Figures in \texttt{results/plots/gap\_reduction\_vs\_gpu\_ntt.png} and
\texttt{competitive\_runtime\_comparison.png}.

% ============================================================
\section{Workload-Adaptive Dispatch}
\label{sec:adaptive}

The gap analysis closed most of the distance to GPU-NTT but did not overturn it everywhere. We
turned the residual \emph{regime structure} into the project's final contribution: a
workload-adaptive NTT system that uses our kernel where it wins and GPU-NTT where it wins.

\textbf{A second, specialized kernel.} To widen the region where we are competitive we added a
\emph{low-latency} kernel (\texttt{tools/low\_latency\_cuda\_ntt/}, our code) specialized at
\emph{compile time} on $\log_2 N$: the stage loop is fully unrolled with compile-time strides and
\verb|__launch_bounds__| pins the thread count to $N/2$, letting the compiler maximize occupancy
in the small-$N$, low-batch latency regime. We then benchmarked \emph{both} our kernels (fused and
low-latency) against GPU-NTT on the same harness, same CUDA-event timing, same modulus, with
forward+inverse round-trip correctness verified for all 56 $(N,\text{batch})$ cells
(\texttt{results/raw/expanded\_gpu\_ntt\_vs\_ours.csv}).

\textbf{The crossover (Figure: \texttt{gpu\_ntt\_gap\_heatmap.png}).} Let
$\text{ratio} = \text{ours\_best} / \text{GPU-NTT}$; $<1$ means we win. Our best kernel wins for
\emph{small transforms} --- $N\le 512$ across every batch (ratio 0.78--0.89), $N=1024$ at batch=1
(0.99), and $N=2048$ at batch $\in\{32,64\}$ (0.96--0.98) --- \textbf{19 of 56 cells}. GPU-NTT
dominates large workloads (up to $4.3\times$ at $N=16384$). Neither system is best everywhere.

\textbf{The selector.} We build an offline-profiled lookup table
(\texttt{results/raw/adaptive\_selector\_table.csv}) mapping each $(N,\text{batch})$ to the fastest
of \{fused, low-latency, GPU-NTT\}; at runtime this is an $O(1)$ table lookup, negligible against a
kernel launch --- the standard autotuning pattern. We evaluate it on four synthetic workloads (200
tasks each), composing total runtime from the \emph{measured} per-call median latencies.

\begin{table}[h]
\centering
\caption{Adaptive selector vs single-system baselines (total workload runtime). Synthetic
workloads, 200 tasks each, composed from measured compute-only latencies.
Data: \texttt{results/raw/mixed\_workload\_comparison.csv}.}
\begin{tabular}{L{3.6cm}R{2.4cm}R{2.2cm}R{2.4cm}}
\toprule
Workload & vs GPU-NTT-only & vs our-only & chose ours/GPU \\
\midrule
A low\_latency\_small  & \textbf{1.14$\times$} & 1.01$\times$ & 148 / 52 \\
B mixed\_realistic     & \textbf{1.02$\times$} & 1.57$\times$ & 50 / 150 \\
C throughput\_large    & 1.00$\times$          & 1.23$\times$ & 0 / 200 \\
D random\_grid         & \textbf{1.04$\times$} & 1.40$\times$ & 80 / 120 \\
\bottomrule
\end{tabular}
\end{table}

The adaptive system is \textbf{never slower} than either single system and is \textbf{strictly
faster than GPU-NTT-only whenever the workload contains small NTTs} (1.14$\times$ on A, 1.04$\times$
on D). On a pure large-$N$ workload (C) it correctly collapses to GPU-NTT ($1.00\times$). Against
our-kernel-only it is up to $1.57\times$ faster by deferring to GPU-NTT where that wins.

\textbf{Honest limitations.} Workload runtimes are summed from measured per-call medians (not a
fresh end-to-end dispatcher run); the selector dispatch cost and inter-task data movement are not
modelled (all timings compute-only). The selector is evaluated on the same grid it was tuned on, so
it is by construction never worse per cell; generalization to unseen $(N,\text{batch})$ (via
nearest-neighbor lookup) is not measured. Our winning region is narrow (small $N$), so the gain
over GPU-NTT-only is modest and appears only when small transforms are present. Single GPU
(RTX~2080~Ti, sm\_75), modulus 998244353, forward NTT. Full discussion:
\texttt{docs/final\_contribution\_strategy.md}.

\subsection{Improvement pass: a warp-shuffle kernel and a real executed runtime}
\label{sec:improvement}

The selector above was evaluated by \emph{summing} measured per-call latencies. Guided by an
options analysis (\texttt{docs/next\_improvement\_strategy.md}), we made two improvements that
strengthen the small-$N$ regime and replace the simulation with a real measurement (Slurm job~93).

\textbf{Warp-shuffle small-$N$ kernel.}
\texttt{tools/low\_latency\_cuda\_ntt/warp\_shuffle\_ntt\_kernels.cu} (our code). In a
decimation-in-time NTT the first five stages have butterfly stride $m<32$, so both partners of
every butterfly lie in one warp. With an element-per-thread layout the partner is $\texttt{lane}
\oplus m$, exchanged in-register via \verb|__shfl_xor_sync| with \textbf{no shared memory and no}
\verb|__syncthreads()|; the cross-warp stages ($m\ge32$) fall back to the proven shared-memory
pair-per-thread path. The kernel is valid for $N\le1024$ (element-per-thread needs
$\text{threads}=N$). Measured, it is the fastest of our kernels on every $N\le1024$ cell, beating
our prior best by $1.10$--$1.18\times$, which \emph{expands and deepens} the region where we beat
GPU-NTT: $N\le512$ improves to ratio $0.61$--$0.74$ (was $0.78$--$0.89$) and $N=1024$ \emph{flips}
from a loss/parity at batch${>}1$ to a $0.85$--$0.95\times$ win. Bit-exact forward+inverse
round-trip verified for all cells. Data: \texttt{results/raw/small\_n\_optimization\_benchmark.csv};
plot \texttt{results/plots/small\_n\_gpu\_ntt\_vs\_ours.png}.

\textbf{A real, executed adaptive runtime.} \texttt{tools/adaptive\_ntt\_runtime/} dispatches each
$(N,\text{batch})$ task to GPU-NTT or our best kernel and times the \emph{whole mixed workload
end-to-end} (CUDA-event wall time, median of 7 runs), rather than summing per-call latencies. The
real run reveals an effect the simulation could not: a \emph{homogeneous} GPU-NTT launch stream
pipelines launches, a \emph{heterogeneous} adaptive stream cannot, and isolated per-call profiling
(one launch+sync per call) overstates per-call cost. An unguarded selector therefore over-routes to
ours; a $10\%$ confidence margin (deviate from GPU-NTT only on robust wins) plus median timing
fixes the small-heavy case and brings broad mixes to within $\sim$1--2\% of GPU-NTT-only.

\begin{table}[h]
\centering
\caption{Real \emph{executed} adaptive runtime vs single-system baselines (median of 7 end-to-end
runs, not summed latencies). Data: \texttt{results/raw/adaptive\_runtime\_benchmark.csv}.}
\begin{tabular}{L{4.0cm}R{2.6cm}R{2.2cm}R{2.4cm}}
\toprule
Workload & vs GPU-NTT-only & vs our-only & routed ours/GPU \\
\midrule
A small-transform-heavy   & \textbf{1.57$\times$} & 1.11$\times$ & 146 / 54 \\
B mixed (large-dominated)  & 0.99$\times$          & 1.84$\times$ & 39 / 161 \\
C throughput-large         & 1.00$\times$          & 1.31$\times$ & 0 / 200 \\
D random grid              & 0.98$\times$          & 1.76$\times$ & 71 / 129 \\
\bottomrule
\end{tabular}
\end{table}

\textbf{Energy and reproducibility.} On the broad mix (workload~D, NVML two-point subtraction) the
adaptive runtime is the \emph{most energy-efficient} policy: $0.098\,\mu$J/NTT and $10.23$M NTT/J
throughput-per-watt, versus GPU-NTT-only's $0.103\,\mu$J/NTT and $9.73$M NTT/J, because our small-$N$
kernels draw $\sim$82\,W vs $\sim$114\,W --- so adaptive trades $\sim$2\% wall-time for $\sim$5\%
better energy efficiency (\texttt{results/raw/adaptive\_energy\_comparison.csv}). Run-to-run
coefficient of variation is $<0.8\%$ everywhere except GPU-NTT on the small-$N$ workload (cv
$\approx20\%$) --- a cold-start instability that is itself the regime our kernels remove
(\texttt{results/raw/adaptive\_reproducibility.csv}; plots
\texttt{adaptive\_runtime\_mixed\_workloads.png}, \texttt{adaptive\_reproducibility\_errorbars.png}).

\textbf{Honest reading.} The adaptive runtime is a decisive win ($1.57\times$) only when small,
low-latency NTTs are a significant fraction of the workload; on large-dominated mixes it is at
parity (within $\sim$1--2\%) but more energy-efficient. We do not force the broad-mix numbers above
parity --- that would mean routing nothing to our kernels.

\subsection{Shared-memory bank-conflict study}
\label{sec:sharedmem}

We then asked whether shared-memory bank conflicts or occupancy are the next bottleneck
(\texttt{docs/shared\_memory\_optimization\_plan.md}, Slurm job~94). The analysis is instructive on
its own: with \texttt{uint32} storage ($4N$ bytes) the per-SM \emph{thread} count (1024) caps blocks
per SM for every $N$, so shared memory is \emph{not} the occupancy limiter; and the
warp-shuffle kernel already performs the bank-conflict-prone early stages (stride ${<}32$) in
registers. So padding can only help the \emph{all-shared} kernels used at $N\ge2048$. We added two
variants (\texttt{tools/shared\_memory\_ntt/}): \textbf{A}, the all-shared kernel with padded indices
$\text{SPAD}(i)=i+\lfloor i/32\rfloor$; and \textbf{B}, the warp-shuffle kernel with its (already
conflict-free) shared phase padded --- a control. Every variant's forward output is bit-for-bit
identical to the verified \texttt{fused} kernel, with round-trip to identity, for all
$N\in\{256,\dots,16384\}$ and batch $\in\{1,8,32,128\}$.

\begin{table}[h]
\centering
\caption{Padding effect on the all-shared kernel (\texttt{low\_latency} $\div$ \texttt{padded};
$>1$ = padding faster). The control (warp-shuffle $\div$ shuffle+pad) was $0.98$--$1.03\times$
everywhere (no effect). $N=16384$ padded exceeds the 64\,KB sm\_75 cap.
Data: \texttt{results/raw/shared\_memory\_optimization\_benchmark.csv}.}
\begin{tabular}{R{1.4cm}R{1.6cm}R{1.6cm}R{1.6cm}R{1.6cm}}
\toprule
$N$ & batch 1 & batch 8 & batch 32 & batch 128 \\
\midrule
256   & 1.01 & 1.01 & 1.00 & 1.01 \\
512   & 0.98 & 1.00 & 1.01 & 1.03 \\
1024  & 1.01 & 1.04 & 1.07 & \textbf{1.23} \\
2048  & 1.11 & 1.07 & 1.12 & \textbf{1.27} \\
4096  & 1.14 & 1.15 & 1.14 & \textbf{1.32} \\
8192  & 1.15 & 1.14 & 1.15 & \textbf{1.29} \\
\bottomrule
\end{tabular}
\end{table}

\textbf{Result.} The control showed no effect (confirming padding only matters where conflicts
exist), but padding the all-shared kernel is \textbf{$1.07$--$1.32\times$ faster at $N\ge2048$}
(neutral at $N\le512$). This \emph{expanded the region where we beat GPU-NTT}: with \texttt{padded}
as the best all-shared kernel, $N=2048$ now wins all tested batches (was losing batch~1--16),
$N=4096$ wins batch~$\ge32$, and $N=8192$ wins batch~128 --- all previously losses. $N=16384$ gets
no benefit (padding does not fit 64\,KB); its gap remains the grid-occupancy limit (one block at
batch=1). An honest correction: we had predicted padding would be neutral, but the all-shared kernel
carried real bank conflicts beyond the early stages, and removing them widened our winning band into
the mid-range. The \texttt{padded} kernel should therefore be added to the adaptive selector for
$N\in\{2048,4096,8192\}$.

\subsection{Integrating \texttt{padded} into the adaptive selector}
\label{sec:padded-integration}

We added \texttt{padded} as a fourth backend to the runtime and built \emph{two} selectors --- OLD
(\{fused, lowlat, shuffle\}) and NEW (\{fused, lowlat, shuffle, \textbf{padded}\}) --- to isolate
its effect (Slurm job~95).

\textbf{Per-call (isolated, compute-only): the winning region expanded.} The NEW selector dispatches
\texttt{padded} for $N=1024$ (batch 64,128), $N=2048$ (all batches), $N=4096$ (batch~$\ge64$), and
$N=8192$ (batch~$\ge64$), flipping $N=4096/8192$ at batch 64/128 from losses to wins and making
$N=2048$ win at every batch --- the winning region grew from ${\sim}24$ to ${\sim}36$ of 56 cells
(\texttt{results/raw/adaptive\_selector\_with\_padded.csv};
\texttt{results/plots/selector\_regions\_with\_padded.png}).

\textbf{Executed mixed workloads: no throughput gain (honest).} End-to-end (median of 7), NEW vs OLD
vs GPU-NTT-only (\texttt{results/raw/mixed\_workload\_with\_padded.csv}):

\begin{table}[h]
\centering
\caption{Integrating \texttt{padded}: executed mixed-workload speedup over GPU-NTT-only. NEW equals
OLD on A/C but slightly regresses B/D, because routing mid-large cells to our kernel forfeits
GPU-NTT's homogeneous-launch pipelining under interleaved execution.}
\begin{tabular}{L{4.0cm}R{2.6cm}R{2.4cm}R{2.6cm}}
\toprule
Workload & NEW (+padded) / GPU & OLD / GPU & padded routes (NEW) \\
\midrule
A small-transform-heavy & $1.23\times$ & $1.23\times$ & 0 \\
B mixed                  & $0.98\times$ & $0.99\times$ & 25 \\
C throughput-large       & $1.00\times$ & $1.00\times$ & 51 \\
D random grid            & $0.97\times$ & $0.98\times$ & 20 \\
\bottomrule
\end{tabular}
\end{table}

So the per-cell isolated win does \emph{not} translate to interleaved throughput. \textbf{Conclusion:}
\texttt{padded} is the correct backend for isolated, latency-bound single transforms at
$N=2048$--$8192$ (a measured per-call win that expanded the winning region, now reflected in the
selector table), whereas for high-throughput interleaved streams the conservative routing is as good.
This is the same offline-profiling-vs-interleaved-execution caveat established for the runtime in
Section~\ref{sec:improvement}.

\subsection{Two selector modes: latency vs throughput}
\label{sec:modes}

The padded result shows that \emph{per-call fastest} selection is not always best for streams. We
therefore expose two policy modes (no new kernels; \texttt{selector\_modes}, job~96):
\textbf{latency mode} (aggressive, margin $0\%$ --- route to the per-call fastest kernel, for
isolated/latency-sensitive calls) and \textbf{throughput mode} (conservative, margin $15\%$ ---
route to our kernel only on a robust win, else GPU-NTT, to avoid harmful backend switching in large
interleaved streams).

\begin{table}[h]
\centering
\caption{Two selector modes vs GPU-NTT-only, executed median of 7 runs
(\texttt{results/raw/adaptive\_selector\_modes.csv}). ``switches'' counts backend changes across the
200-task stream. Throughput mode is $\ge$ latency mode on every workload.}
\begin{tabular}{L{3.8cm}R{2.2cm}R{2.4cm}R{2.6cm}}
\toprule
Workload & latency / GPU & throughput / GPU & switches (lat / thr) \\
\midrule
A small-transform-heavy & $1.15\times$ & $\mathbf{1.34\times}$ & 0 / 81 \\
B mixed                  & $0.91\times$ & $0.99\times$ & 132 / 70 \\
C throughput-large       & $0.98\times$ & $0.99\times$ & 104 / 0 \\
D random grid            & $0.93\times$ & $0.98\times$ & 132 / 97 \\
\bottomrule
\end{tabular}
\end{table}

\textbf{Latency mode} over-routes marginal mid/large cells to our kernel and thrashes backends (132
switches on B/D), forfeiting GPU-NTT's launch pipelining and regressing those streams to
$0.91$--$0.93\times$. \textbf{Throughput mode} delivers the large small-NTT win (A, $1.34\times$ ---
even beating latency mode there by offloading the marginal $N{=}1024$ cells to GPU-NTT) and matches
GPU-NTT-only within ${\sim}1$--$2\%$ (measurement noise) on the large-dominated streams. The correct
policy is therefore workload-aware: latency mode for isolated calls, throughput mode for streams.

\subsection{Architecture pass: multi-block 4-step NTT for low-batch large $N$}
\label{sec:multiblock}

A measured \texttt{ptxas\,-v} audit of every kernel (\texttt{results/raw/kernel\_resource\_report.csv})
ruled out the usual suspects: \textbf{zero register spills anywhere}, and no kernel is
register-bound (18--45 registers). It isolated a single remaining architectural limiter ---
\textbf{grid-level under-occupancy at low batch}. For $N\ge 2048$ the all-shared kernels launch
1024-thread blocks, exhausting a Turing SM's entire thread budget ($1024$ threads/SM) so that only
\emph{one} block resides per SM; combined with the one-block-per-transform design, at batch=1 a
single block runs on \textbf{1 of the GPU's 68 SMs ($1.5\%$ device utilization)} as a $14$-stage
serial chain. This exactly explains the measured low-batch large-$N$ gap, which grows monotonically
from $1.27\times$ at $N{=}16384$, batch=128 to $4.31\times$ at batch=1.

The fix is a \textbf{multi-block 4-step (Cooley--Tukey) decomposition}. Factoring $N=N_1 N_2$ (e.g.
$16384=128\times128$) with the index split $n=n_1+N_1 n_2$, $k=N_2 k_1 + k_2$ gives
$X[N_2 k_1 + k_2] = \sum_{n_1} w_{N_1}^{n_1 k_1}\, w_N^{n_1 k_2}\,\big[\sum_{n_2} x[n_1+N_1 n_2]\,
w_{N_2}^{n_2 k_2}\big]$. Kernel \texttt{mb\_step1} computes the $N_1$ independent inner $N_2$-point
sub-NTTs plus the $w_N^{n_1 k_2}$ twiddle; \texttt{mb\_step2} computes the $N_2$ independent outer
$N_1$-point sub-NTTs. One transform now launches $N_1$ (then $N_2$) blocks of tiny $\le128$-point
sub-NTTs on $\le 64$ threads each (up to $16$ blocks/SM), \textbf{saturating all 68 SMs at batch=1}
and halving the per-block serial chain ($14\to7$ stages). The output index $N_2 k_1 + k_2$ is a
bijection of $[0,N)$, so results land in natural order with no extra permutation. Correctness is
bit-exact versus the CPU DIT reference and a forward$+$inverse round-trip on all 32 cells
($N\in\{2048,\dots,16384\}\times$ batch $\in\{1,\dots,128\}$); SLURM jobs 105/106.

\begin{table}[h]
\centering
\caption{Multi-block 4-step NTT, compute-only (CUDA events, median of 30), modulus 998244353.
Gap $=$ ours $\div$ GPU-NTT (lower is better; $<1$ we win).}
\begin{tabular}{rrrrrcc}
\toprule
$N$ & batch & prior best (ms) & \textbf{multi-block (ms)} & GPU-NTT (ms) & prior gap & \textbf{new gap} \\
\midrule
4096  & 1 & 0.0173 & \textbf{0.0116} & 0.0120 & $1.50\times$ & \textbf{0.97$\times$ (win)} \\
8192  & 1 & 0.0278 & \textbf{0.0124} & 0.0117 & $2.50\times$ & \textbf{1.06$\times$} \\
16384 & 1 & 0.0512 & \textbf{0.0142} & 0.0122 & $4.31\times$ & \textbf{1.16$\times$} \\
16384 & 8 & 0.0517 & \textbf{0.0311} & 0.0164 & $3.20\times$ & \textbf{1.90$\times$} \\
\bottomrule
\end{tabular}
\end{table}

The \textbf{$N{=}16384$, batch=1 gap --- the project's single worst case --- falls from
$4.31\times$ to $1.16\times$} ($3.60\times$ faster than our prior best), and $N{=}4096$, batch=1
flips to a \emph{win} over GPU-NTT. The design adds a second kernel launch and ${\sim}4$ DRAM passes
(versus $2$): free when the GPU is $98\%$ idle, but a net cost once batch fills the grid, so the
single-block fused kernel retakes the lead at batch $\ge 16$--$32$ (measured crossover). Multi-block
is therefore adopted as a \emph{new selector backend used only in its winning corner}
($N\ge4096$, low batch); no other regime regresses. This confirms the bottleneck matrix's
prediction and mirrors the multi-block strategy GPU-NTT itself uses for large $N$. Full analysis:
\texttt{docs/architecture\_optimization\_strategy.md}.

\paragraph{Stream-level integration (executed adaptive runtime).} We added multi-block as a
\emph{fifth} backend to the real dispatcher (\texttt{tools/adaptive\_ntt\_runtime/}) and re-measured
end-to-end (SLURM job 107, median of 7 over workloads A--E; all 56 profile cells forward$+$round-trip
correct). Multi-block is the fastest of \emph{our} kernels for $N{=}4096$ ($b{\le}8$), $N{=}8192$
($b{\le}8$), $N{=}16384$ ($b{\le}16$); single-block wins above. Versus GPU-NTT it reaches only
\emph{parity} there (ratio ${\approx}1.0$--$1.16$, below the $15\%$ throughput margin), so the
conservative selector keeps GPU-NTT --- the correct call, since GPU-NTT pipelines a homogeneous
large-$N$ launch stream that a heterogeneous dispatch cannot reclaim from a per-call tie. The
executed result is therefore \textbf{parity with no regression} on every workload (A small-heavy
$1.56\times$ unchanged, B $0.99\times$, C $1.03\times$, D $0.98\times$, E $1.00\times$ vs
GPU-NTT-only), while the \emph{standalone} our-only system improves dramatically: on the new large-$N$
low-batch workload E our-only drops from $4.16$ to $2.08$\,ms ($\mathbf{1.995\times}$, cv $0.05\%$),
and on the mixed/random workloads by $1.20$--$1.22\times$. The naïve latency mode that force-routes
the at-parity cells regresses B/D to $0.92$--$0.93\times$ (132 backend switches), confirming the
per-call tie does not survive as a stream win. \textbf{Honest summary:} the multi-block kernel fixes
the low-batch grid under-occupancy so our standalone system is competitive across the whole
$(N,\text{batch})$ space instead of $3$--$4\times$ behind in one corner, and the adaptive runtime
chooses single-block or GPU-NTT once the grid is saturated --- so the integrated system never
regresses below GPU-NTT-only while preserving the $1.56\times$ small-transform win.

\begin{enumerate}
  \item \textbf{Stage fusion in Triton}: Once Triton 3.3+ adds \texttt{tl.gather},
    implement multi-stage kernel fusion (Merge NTT style) to eliminate $\log_2
    N - \log_2 N / d$ DRAM passes for merge depth $d$.
  \item \textbf{Mixed-radix NTT}: Radix-4 or radix-8 butterflies halve or quarter the
    stage count, reducing both launches and DRAM traffic proportionally.
  \item \textbf{Profile with Nsight Compute}: Measure actual DRAM bytes, SM occupancy,
    and arithmetic throughput to validate the traffic model.
  \item \textbf{GPU-NTT comparison and gap closure}: \emph{Done} — built from source and
    benchmarked at $998244353$ (Section~\ref{sec:external}); then a fused uint32 kernel
    (Section~\ref{sec:competitive}) closed the gap to 1.1--1.3$\times$ at batch=128, and a
    multi-block 4-step kernel (Section~\ref{sec:multiblock}) closed the batch=1 occupancy limit
    ($N{=}16384$ b=1: $4.31\times\!\to\!1.16\times$). Remaining work: an explicit Barrett-$\mu$
    butterfly. ICICLE's GPU NTT remains unavailable (closed-source backend).
  \item \textbf{Larger $N$}: Test $N$ up to $2^{23}$ (8.4\,M), the maximum for
    modulus 998244353, to match cryptographic CKKS/BFV parameters.
\end{enumerate}

% ============================================================
\section{Post-Quantum Cryptography: the Kyber NTT}
\label{sec:pqc}

Everything so far used the NTT-friendly modulus $q=998244353$. To close the loop on the project's
stated motivation (lattice PQC), we adapted and characterized the pipeline on a \emph{real}
post-quantum parameter set --- \textbf{Kyber / ML-KEM: $n=256$, $q=3329$}. We do \emph{not} claim to
beat reference (AVX2) Kyber; the contribution is (i) a correct GPU Kyber poly-mul, (ii) a roofline
re-analysis of the small-modulus regime, and (iii) a direct test of our earlier Montgomery prediction.

\textbf{Negacyclic NTT and why Kyber needs an \emph{incomplete} one.} PQC multiplies in
$\mathbb{Z}_q[x]/(x^n+1)$ (negacyclic). At $q=998244353$ this is done by $\psi$-twisting ($\psi$ a
primitive $2n$-th root), reusing our fused kernel unchanged (verified bit-exact, job 97). Kyber is
different: $q-1 = 2^8\cdot 13$, so its 2-adicity is only $2^8=256$ and \textbf{no primitive $512$-th
root exists} --- a full negacyclic-256 transform is impossible. We therefore implement the real Kyber
\textbf{incomplete NTT}: 7 layers ($\zeta=17$) mapping $\mathbb{Z}_q[x]/(x^{256}+1)$ to 128 rings
$\mathbb{Z}_q[x]/(x^2-\gamma_i)$, with ``pointwise'' replaced by 128 base-case linear-polynomial
products. The whole poly-mul is fused in shared memory (1\,KB/poly), verified bit-exact against an
independent CPU schoolbook reference (20 pairs + batch-32, job 99).

\textbf{Step 3 --- roofline re-analysis.} The large modulus was bandwidth-bound (80--84\% of peak BW).
The fully-resident Kyber kernel is the opposite: 100\% theoretical occupancy, arithmetic intensity
$1.167$ modmul/byte, and achieved DRAM bandwidth peaking at only \textbf{26.6\% of peak} (Figure
\texttt{kyber\_roofline.png}). It is \textbf{not bandwidth-bound} --- nor compute-bound (191 Gmodmul/s,
$\sim$6\% of the integer ceiling) --- but \textbf{latency/occupancy-bound}: 14 \verb|__syncthreads| per
poly-mul and few blocks at small batch.

\textbf{Step 4 --- the Montgomery prediction is refuted (with a reason).} Our earlier finding was that
Montgomery hurts when bandwidth-bound; the hypothesis was that it should help once compute-bound and
fully resident (Kyber). Measured, Montgomery is \textbf{1.18--1.55$\times$ slower} than Barrett at
every batch and uses more energy ($3.03$ vs $2.31\,\mu$J/poly-mul).

\begin{table}[h]
\centering
\caption{Kyber poly-mul: Montgomery vs Barrett (RTX~2080~Ti, compute-only). Montgomery loses at every
batch. Data: \texttt{results/raw/kyber\_bench.csv}.}
\begin{tabular}{R{1.6cm}R{2.0cm}R{2.2cm}R{2.4cm}}
\toprule
batch & Barrett ($\mu$s) & Montgomery ($\mu$s) & Mont/Barrett \\
\midrule
1    & 4.99  & 7.76  & \textbf{1.55$\times$} \\
128  & 6.11  & 8.58  & 1.40$\times$ \\
512  & 12.6  & 15.0  & 1.19$\times$ \\
2048 & 38.4  & 45.4  & \textbf{1.18$\times$} \\
\bottomrule
\end{tabular}
\end{table}

The reason is precise: Montgomery's advantage is avoiding an \emph{expensive} reduction, which is a
property of \emph{modulus width}, not of being compute-bound. At $q=3329$ a product is
$<(q-1)^2<2^{24}$, fits in one 32-bit word, and Barrett is two 32-bit multiplies plus a conditional
subtract --- already minimal. Montgomery gives no cheaper reduction and adds domain conversions, so it
is a net loss. ``Compute-bound $\Rightarrow$ Montgomery wins'' is too coarse; \textbf{Barrett is the
right choice for small ($\le$16-bit) PQC moduli}, while the earlier 30-bit analysis still stands.

\textbf{Step 5 --- batched throughput and energy.} Compute-only on the RTX~2080~Ti, batched Kyber
poly-mul reaches \textbf{53.4\,M poly-mul/s} and \textbf{128\,M forward-NTT/s} at batch 2048 ---
\textbf{456$\times$} a single-thread CPU incomplete-NTT baseline --- at \textbf{2.31\,$\mu$J/poly-mul}
(NVML, two-point). GPU-NTT and ICICLE are \emph{not} benchmarked at $q=3329$ (their small-field paths
target babybear / large NTT-friendly primes, not 3329); the in-scope baseline is our CPU
implementation. Full numbers, plots, and honest limitations: \texttt{docs/pqc\_kyber\_analysis.md}.

\subsection{Fused Kyber matrix--vector core: the right unit of work}
After the isolated poly-mul reached its architectural optimum (removing barriers cost registers and
occupancy, giving no net gain), we changed the \emph{unit of work}. Real Kyber computes a
\textbf{matrix--vector product} $t = A\circ s$, where $A$ is a $k\times k$ matrix of NTT-domain
polynomials and $s$ a $k$-vector. This exposes two reuses the isolated poly-mul cannot: (i)
\textbf{NTT reuse} --- each $s_j$ is forward-transformed \emph{once} and reused across all $k$ rows;
(ii) \textbf{NTT-domain accumulation} --- the $k$ base-muls of a row are summed in the NTT domain
before a \emph{single} inverse NTT. One matrix--vector product then needs \textbf{$2k$ NTT-transforms}
instead of the ${\sim}3k^2$ of $k^2$ isolated poly-muls. The kernel \texttt{kyber\_matvec<K>} (one block
per instance, 128 threads; 44/46/52 registers, 0 spills, 100\% occupancy) is verified bit-exact against
an independent CPU matrix--vector reference (INTT $+$ negacyclic schoolbook) for $k=2,3,4$ over 205
instances each (random $+$ edge), all outputs canonical in $[0,q)$.

\begin{center}
\begin{tabular}{lcccc}
\hline
$k$ (Kyber) & vs $k^2$ isolated (batch 2048) & eff.\ poly-mul/s & energy/poly-mul & poly-mul/J \\
\hline
2 (512)  & $2.39\times$ & 148\,M & --- & --- \\
3 (768)  & $3.48\times$ & 220\,M & --- & --- \\
4 (1024) & $\mathbf{4.44\times}$ & $\mathbf{283}$\,M & $\mathbf{0.68\,\mu}$J ($-78\%$) & $4.6\times$ \\
\hline
\end{tabular}
\end{center}

The speedup grows with $k$ (more NTT reuse) and with batch (the grid fills), and it is the operation
real server-side Kyber actually performs. Effective Kyber throughput rises from the prior 56\,M
poly-mul/s (isolated) to \textbf{283\,M} ($k{=}4$, batch 2048), at $4.6\times$ better energy per
operation. We implement only the arithmetic core $A\circ s$, not full ML-KEM.

\textbf{Two follow-ups, honestly reported.} (a) \emph{Shoup/Harvey butterflies} for the large-modulus
NTT ($q=998244353$; precomputed $w'=\lfloor w\,2^{32}/q\rfloor$, one \texttt{\_\_umulhi}, overflow-proven)
give a scoped $\sim\!1.04$--$1.06\times$ over our own baseline at $N\le 4096$ and large-$N$ low batch,
bit-exact at every $(N,\text{batch})$; interleaving $w$ and $w'$ into one \texttt{uint2} load removes an
initial large-$N$ regression, proving that regression was doubled twiddle-table DRAM traffic --- but the
gain is below $1.20\times$. (b) The same Shoup lever \emph{inside} the matvec NTT is a clean negative
($0.95$--$1.05\times$): the matvec is \texttt{\_\_syncthreads}/shared-memory-latency bound, not
arithmetic-bound, confirming its win is algorithmic. Full record:
\texttt{docs/autonomous\_optimization\_results.md}.

\subsection{End-to-end batched Kyber arithmetic pipeline ($t = A\circ s + e$)}
We integrated the fused core into a realistic \textbf{end-to-end} pipeline with \textbf{int16 I/O} and
full layout conversion --- the Kyber matrix-vector-plus-error operation $t = A\circ s + e$ ($\hat A$ in
the NTT domain, $s,e$ coefficient-domain $k$-vectors, output canonical) --- measured as \textbf{full
executed wall time}. Scope is the arithmetic core only (no sampling/hashing/compression/serialization;
not full ML-KEM). Three variants: \textbf{A} isolated (\texttt{kyber\_polymul}, $3k^2$ transforms),
\textbf{B} fused multi-kernel (\texttt{kyber\_matvec}, $2k$), \textbf{C} fused-integrated
(\texttt{kyber\_matvec\_e2e}: int16 read $+$ matvec $+$ error-add $+$ reduction $+$ int16 store in one
launch). All three are \textbf{bit-exact vs an independent CPU $A\circ s+e$ reference} ($k=2,3,4$, 205
instances each, 0 out-of-range).

\begin{center}
\begin{tabular}{lcccc}
\hline
$k$ & full-pipeline C/A (batch 2048) & energy/op (C) & energy/op (A) & eff.\ eq-products/s (C) \\
\hline
2 & $4.22\times$ & --- & --- & 151\,M \\
3 & $5.97\times$ & --- & --- & 227\,M \\
4 & $\mathbf{7.48\times}$ & $\mathbf{10.3\,\mu}$J & $79.8\,\mu$J & $\mathbf{296}$\,M \\
\hline
\end{tabular}
\end{center}

The full-pipeline speedup ($\le 7.48\times$) \emph{exceeds} the kernel-only $4.44\times$ because the
isolated pipeline also pays conversion, gather, and accumulate over $3k^2$ transforms; energy per
complete operation drops $7.7\times$ ($79.8\!\to\!10.3\,\mu$J, NVML). A stage breakdown shows the
int16$\to$u32 conversion pass alone ($0.127$\,ms) exceeds the whole fused kernel ($0.111$\,ms), and $C$
beats even the raw u32 matvec because int16 I/O halves the matrix DRAM read. CUDA Graphs help the
multi-kernel path at small batch (up to $1.56\times$); mixed-$k$ grouped dispatch avoids a per-instance
launch explosion. A Phase-9 attempt to add the Shoup-v2 kernel to the executed large-modulus selector is
an honest negative (routed 0 times; the warp-shuffle / padded kernels already dominate its cells). Full
record: \texttt{docs/kyber\_end\_to\_end\_arithmetic\_plan.md}.

% ============================================================
\section{Final Authoritative Results (2026-07-06)}

This section consolidates the final same-harness evidence (job 139, srvr-2, adaptive runtime
\texttt{shoup\_v2}, $q=998244353$, CUDA events, WARM=3, REP=9, compute-only). Raw data:
\texttt{results/raw/final\_ntt\_performance\_map.csv} and \texttt{final\_ntt\_vs\_gpu\_ntt.csv};
architecture summary: \texttt{docs/final\_ntt\_architecture\_summary.md}.

\subsection{Final adaptive CUDA system and the GPU-NTT comparison (74-cell map)}
The final measured map covers \textbf{74 cells}: $N\in\{256,\dots,16384\}$ at batch $1\dots128$,
plus an extended-batch sweep to $8192$ for $N\le1024$ (authoritative run: job 153, srvr-2, all cells in
one program with the selected backend \emph{executed} at REP=100 and values+ordering gated). Under the
raw native-order contract our kernels beat GPU-NTT (${>}1.03\times$) in \textbf{48 cells} --- maxima
\textbf{1.93$\times$} ($N{=}256$), 1.62$\times$ (512), 1.48$\times$ (1024), 1.38$\times$ (2048),
1.34$\times$ (4096), 1.18$\times$ (8192) --- tie within 3\% in \textbf{4}, and GPU-NTT is faster in
\textbf{22} (the DRAM-saturated region, where its per-stage Merge kernels stream better). A second run
of the same grid (job 151) gave 49/7/18 with max 2.14$\times$: boundary cells flip within $\pm$3\%
noise. Selector execution: regret (executed p50 $\div$ best measured) median 0.998, worst 1.051;
table-lookup overhead 1.25\,ns. Energy (NVML counter, mixed workload D): adaptive 0.192\,$\mu$J/NTT
at 255\,W.
Small transforms are won by the warp-shuffle+Shoup kernel (22 registers, 0 spills, 5 sync-free
shuffle stages, 100\% theoretical occupancy); the extended sweep shows this win \emph{persists to
batch $\approx$2048--4096 at $N{=}256$} before saturation. Mid sizes are won by the bank-conflict-free
padded(+Shoup) kernel; large-$N$ low-batch cells by the multi-block 4-step kernel (128+ blocks fill
all 68 SMs; argmin recovers up to 1.15$\times$ at $N{=}4096$ $B{=}2$). The argmin selector
(one-constant change from the previous 10\%-margin rule) routes 44/56 base cells to our kernels vs
37 before, recovering 7 boundary cells with no measured stream regression; with GPU-NTT as a fallback
backend the runtime is \emph{never worse than GPU-NTT} on any measured cell.

\textbf{Correctness and convention.} All backends pass edge-input gates (zero, one-hot, maximum
residue $q{-}1$, random) plus inverse round trips at every size. The output-ordering convention was
formally verified: our kernels emit the \emph{natural-order} cyclic NTT (bit-exact with the direct
DFT at $\omega=3^{(P-1)/N}$; bit-reversal folded into the load), while GPU-NTT emits the same values
in \emph{bit-reversed} order --- identical transforms, so the comparison charges no reordering cost
to either side.

\begin{figure}[htbp]\centering
\includegraphics[width=0.86\textwidth]{../results/plots/final_ntt_vs_gpu_ntt_heatmap.png}
\caption{Best-of-ours $\div$ GPU-NTT per $(N,\text{batch})$ cell. Green: our kernels faster; the
adaptive runtime selects the per-cell winner (never worse than GPU-NTT).}
\end{figure}

\subsection{Representative throughput (execution-verified, raw native order)}
\begin{center}\small
\begin{tabular}{llllrrrr}
\toprule
workload & $N$ & batch & backend & $\mu$s/NTT & NTT/s & eff.\,GB/s & vs GPU-NTT \\
\midrule
small-$N$ low batch & 256 & 1 & shoup\_shuf & 4.58 & 0.219M & 0.4 (0.1\%) & 1.93$\times$ \\
small-$N$ pre-saturation & 256 & 2048 & shoup\_shuf & 0.010 & 96.2M & 197 (32\%) & 1.08$\times$ \\
small-$N$ saturation & 256 & 8192 & gpu\_ntt & 0.009 & 117.2M & 240 (39\%) & 0.95 \\
medium-$N$ & 2048 & 32 & shoup\_pad & 0.280 & 3.57M & 59 (9.5\%) & 1.33$\times$ \\
large-$N$ low batch & 16384 & 1 & gpu\_ntt & 10.24 & 98K & 13 (2.1\%) & 0.97 \\
large-$N$ high batch & 16384 & 128 & gpu\_ntt & 0.701 & 1.43M & 187 (30\%) & 0.82 \\
\bottomrule
\end{tabular}
\end{center}
Effective GB/s is a traffic estimate (8\,B/coefficient $\div$ measured runtime), not a profiler
counter. Our-kernel peak: 110.9M NTT/s (warp\_shuffle\_shoup, $N{=}256$, batch 8192, 0.0739\,ms).

\subsection{Optimization progression}
Figure in \texttt{results/plots/final\_ntt\_optimization\_progression.png} and
\texttt{results/raw/final\_ntt\_optimization\_progression.csv} trace each stage
(CPU $\to$ C++ $\to$ PyTorch $\to$ Triton radix-2 $\to$ radix-4 $\to$ fused CUDA $\to$ warp-shuffle
$\to$ padded $\to$ multi-block $\to$ Shoup $\to$ adaptive), with the bottleneck each targeted and its
measured improvement in its regime (timing harness noted; the kernels are regime-specialized by design).

\subsection{Kyber as an NTT case study}
Applying the same NTT principles to Kyber's incomplete seven-layer NTT ($q=3329$, $n=256$), at $k=4$,
batch~2048: the isolated-transform baseline (0.828\,ms) improves to a fused matrix-vector kernel with
shared-memory transform reuse (0.276\,ms, 3.0$\times$), then a single-kernel fused-integrated kernel with
int16 I/O (0.111\,ms, 2.48$\times$), then NTT-domain accumulation that drops the final INTT (0.067\,ms,
1.65$\times$) --- \textbf{12.3$\times$ total latency and $\sim$12.7$\times$ energy}
(79.8$\to$6.3\,$\mu$J/op), each step attributed to a specific NTT optimization principle and verified
bit-exact against an independent CPU schoolbook reference for $k=2,3,4$. Data:
\texttt{results/raw/kyber\_ntt\_contribution\_breakdown.csv}.

\begin{figure}[htbp]\centering
\includegraphics[width=0.9\textwidth]{../results/plots/kyber_ntt_contribution_breakdown.png}
\caption{Kyber arithmetic: each step reuses a GPU NTT principle (shared-memory transform reuse,
single-kernel fusion, NTT-domain accumulation). The improvement is NTT-attributed, not from system
batching or transfer routing.}
\end{figure}

\subsection{Future work and system extensions}
The following are working, correctness-verified extensions built while validating the NTT result, but are
deliberately out of the core NTT narrative: a full-scope CPU ML-KEM baseline (pq-crystals AVX2) for
falsifiable positioning; an on-GPU FIPS-203 seed generator with a measured transfer-vs-seed crossover
router; server window-grouping throughput/tail analysis; and CUDA-graph small-batch experiments. They
demonstrate the NTT kernels inside a realistic PQC workload but do not change the NTT contribution.

\section{Conclusion}

We implemented a GPU-accelerated NTT pipeline progressing from Python CPU to a custom
Triton GPU kernel, with progressively deeper analysis of the bottlenecks at each stage.

\textbf{Optimization 1 (PyTorch level)}: Precomputed cached twiddle factors and
batch-parallel tensor operations give $\sim$5$\times$ improvement over a simple GPU
implementation, with cache break-even in under one call.

\textbf{Optimization 2 (kernel level)}: A custom Triton butterfly kernel with 1 launch
per stage (vs 7 PyTorch ops) and LLVM-generated Barrett reduction (via \texttt{tl.constexpr}
MOD) achieves \textbf{3.20--4.65$\times$ additional speedup} over the PyTorch baseline.
For large problem sizes, the Triton kernel reaches 82\% of peak memory bandwidth (504\,GB/s)
and the speedup is primarily explained by 4.27$\times$ less DRAM traffic.

Against the C++ CPU baseline, Triton achieves up to \textbf{58.6$\times$ speedup} at
batch=128, $N=16384$, reflecting both GPU parallelism and bandwidth advantages.

\textbf{Optimization 3 (closing the gap to the state of the art)}: Benchmarking the open-source
GPU-NTT library on our own GPU first showed it 8--30$\times$ faster. A gap analysis attributed
this to int64 storage, per-stage kernel launches, and a separate bit-reversal pass. We
implemented a single-kernel fully-fused \texttt{uint32} NTT (bit-reversal folded into the load)
and, measuring it identically to GPU-NTT, \textbf{reduced the gap to 1.1--1.3$\times$ at batch=128
across all $N$ and to parity/advantage at small $N$} --- a 6--32$\times$ improvement over our prior
best. This demonstrates concretely which GPU NTT optimizations matter most: launch-count
reduction via stage fusion and halved DRAM traffic via 32-bit storage. The contribution is the
measured analysis plus the optimization that closes most of the gap, with the residual large-$N$
batch=1 gap honestly attributed to single-block occupancy.

\textbf{Optimization 4 (a workload-adaptive system)}: Rather than report a single global winner,
we added two specialized kernels --- a compile-time-specialized low-latency kernel and a
warp-shuffle kernel (Section~\ref{sec:improvement}) that drops shared memory and
\verb|__syncthreads()| for the within-warp early stages --- establishing a measured regime
crossover (our kernels win for small transforms; GPU-NTT wins for large). The warp-shuffle kernel
\emph{expanded} that winning region: $N\le512$ deepened to $0.61$--$0.74\times$ and $N=1024$ flipped
to a $0.85$--$0.95\times$ win over GPU-NTT. We then built a \emph{real, executed} hybrid runtime
(not a latency simulation) that dispatches per $(N,\text{batch})$ and times whole mixed workloads
end-to-end. It is \textbf{$1.57\times$ faster than GPU-NTT-only on small-transform-heavy workloads},
at parity (within $\sim$1--2\%) on large-dominated mixes, and the most energy-efficient policy on a
broad mix. A bank-conflict-padded kernel (Section~\ref{sec:sharedmem}) then extended our per-call
wins to $N{=}2048$--$8192$, and we exposed the runtime as \textbf{two workload-aware modes}
(Section~\ref{sec:modes}): a \emph{latency mode} (per-call fastest, for isolated calls) and a
\emph{throughput mode} (conservative, robust-win-only, for interleaved streams). The final claim is
therefore not ``we beat GPU-NTT'' but the more honest and more useful one: \emph{our project
contributes a workload-aware adaptive NTT runtime with two modes --- a latency mode that exploits our
low-latency/padded kernels where they win, and a throughput mode that avoids harmful backend
switching and falls back to GPU-NTT for large homogeneous streams --- delivering up to
$1.34$--$1.57\times$ over GPU-NTT-only on small-transform-heavy workloads while matching it on
large ones.}

All unit tests pass with bit-exact cross-validation between implementations, including the new
competitive kernel (round-trip, direct-DFT, and polynomial-multiplication checks).

% ============================================================
\begin{thebibliography}{9}

\bibitem{gpuntt}
Ali Şah Özcan, Arsalan Javeed, and Erkay Savaş.
\textit{High-Performance Number Theoretic Transform on GPU Through radix2-CT and 4-Step Algorithms.}
IEEE Access, vol.~13, pp.~87862--87883, 2025.
DOI: \url{10.1109/ACCESS.2025.3570024}.
Also available as: ePrint 2023/1410.
Repository: \url{https://github.com/Alisah-Ozcan/GPU-NTT} (MIT License).

\bibitem{icicle}
Ingonyama.
\textit{ICICLE: A GPU-Accelerated Cryptography Library for Zero-Knowledge Proofs.}
\url{https://github.com/ingonyama-zk/icicle} (MIT License, frontend).

\end{thebibliography}

\end{document}
